\documentclass[book.tex]{subfiles}
%\usepackage{subfiles}
\begin{document}

%\pagestyle{empty}
\vspace*{36 pt}
\part{Appendix: Mathematical Foundations}
\label{part:math}

This part provides a brief introduction to a variety of different mathematical topics. If you're working through this book for the first time, you may prefer to skim sections in this part, and then come back and refer to sections as you encounter the mathematics in later chapters.

Each chapter has a number of associated exercises. If you're new to a topic, work through the exercises to gain some practice with the concept. The only way to internalize mathematics is practice. If you're new to formal mathematics, don't worry about the formalism; instead, try to build intuition. You may find yourself coming back to this part multiple times as you work your way through the rest of the book.
\newpage
\subfile{chapters/part1: Mathematical Foundations/math_basics}
\subfile{chapters/part1: Mathematical Foundations/calculus}
\subfile{chapters/part1: Mathematical Foundations/vectors_scalars_tensors}
\subfile{chapters/part1: Mathematical Foundations/ordinary_differential_equations}
\subfile{chapters/part1: Mathematical Foundations/sturm_liouville_theory}
\subfile{chapters/part1: Mathematical Foundations/probability_and_statistics}
\subfile{chapters/part1: Mathematical Foundations/information_theory}
\subfile{chapters/part1: Mathematical Foundations/generating_functions}
\subfile{chapters/part1: Mathematical Foundations/matrix_calculus}
\subfile{chapters/part1: Mathematical Foundations/distributions}
\subfile{chapters/part1: Mathematical Foundations/partial_differential_equations}
\subfile{chapters/part1: Mathematical Foundations/greens_equations}
\subfile{chapters/part1: Mathematical Foundations/hilbert_spaces}
\subfile{chapters/part1: Mathematical Foundations/optimization}
\subfile{chapters/part1: Mathematical Foundations/special_functions}
\subfile{chapters/part1: Mathematical Foundations/complex_analysis}
\subfile{chapters/part1: Mathematical Foundations/manifolds}
\subfile{chapters/part3: Physics/general_relativity/tensor_analysis}
\subfile{chapters/part1: Mathematical Foundations/covariant_derivative_connection}
\subfile{chapters/part1: Mathematical Foundations/symplectic_manifold}
\subfile{chapters/part1: Mathematical Foundations/lie_group_and_lie_algebra}
\subfile{chapters/part1: Mathematical Foundations/lorentz_group}
\subfile{chapters/part1: Mathematical Foundations/spinors}
% Conv as an operator and link back to 1c
% TODO \subfile{chapters/part1: Mathematical Foundations/dynamical_systems}
\subfile{chapters/part1: Mathematical Foundations/mathematical_transforms}
\subfile{chapters/part1: Mathematical Foundations/fiber_bundles}




\end{document}